% =====================================================================
\section{Computer Vision: How a UAV Sees the Rovers}\label{sec:vision}
\pipestrip{9}
Each UAV carries a thermal camera on its gimbal. The autonomy never reads a rover's position: everything it knows comes from a $320\times240$ image of apparent temperature. Figure~\ref{fig:cvoverview} follows one frame to a confirmed, fused, predicted track.
\begin{figure}[h]
\centering
\begin{tikzpicture}[node distance=3mm and 4mm, every node/.style={font=\scriptsize},
  box/.style={draw,rounded corners=2pt,minimum height=8mm,align=center,fill=blue!6},arr/.style={-{Latex[length=1.6mm]},thick}]
  \node[box] (a) {heat render\\$320{\times}240$};
  \node[box,right=of a] (b) {threshold\\Otsu + floor};
  \node[box,right=of b] (c) {despeckle};
  \node[box,right=of c] (d) {labelling\\union-find};
  \node[box,right=of d] (e) {moments,\\filter};
  \node[box,below=of a] (f) {pixel ray\\$\to$ ground point};
  \node[box,right=of f] (g) {gating\\$\chi^2=9.21$};
  \node[box,right=of g] (h) {Hungarian};
  \node[box,right=of h] (i) {Kalman\\update};
  \node[box,right=of i] (j) {two-view\\triangulation};
  \foreach \x/\y in {a/b,b/c,c/d,d/e,f/g,g/h,h/i,i/j} \draw[arr] (\x)--(\y);
  \draw[arr] (e.south) -- ++(0,-3mm) -| (f.north);
\end{tikzpicture}
\caption{The vision chain, per UAV frame (top) and shared by the squad (bottom).}\label{fig:cvoverview}
\end{figure}

\subsection{The thermal image}
\autofig{diagrams}{G2_cv_pipeline}
The scene is drawn again from the gimbal pose into a $320\times240$ target ($40^\circ$ vertical field of view) with the \emph{raw heat} output instead of colour. Each material has an intrinsic temperature $h_0$ (ground $0.15$, bunker $0.34$, a rover's motor box $0.97$, a charred wreck $0.95$ cooling
as $0.3+0.65e^{-t/40\,\text{s}}$) plus a little absorbed light and ground noise; the result is read back as 8-bit values (heat $\times255$). Rovers are the only ground-level parts hotter than $0.9$; static props stay below the blob threshold (T41: hottest prop $0.74$), as do tyre footprints ($\le0.51$) and the wreck.
\begin{figure}[h]
\centering
\begin{subfigure}{0.49\linewidth}\piclab{cv_thermal}{\linewidth}\caption{raw heat image (UAV 0)}\end{subfigure}
\begin{subfigure}{0.49\linewidth}\piclab{cv_threshold}{\linewidth}\caption{after Otsu + floor threshold}\end{subfigure}\\[1mm]
\begin{subfigure}{0.49\linewidth}\piclab{cv_despeckle}{\linewidth}\caption{after despeckle}\end{subfigure}
\begin{subfigure}{0.49\linewidth}\piclab{cv_labels}{\linewidth}\caption{labelled blobs, centroids, orientation}\end{subfigure}
\caption{The image stage on a real frame, one picture per step.}\label{fig:cvstages}
\end{figure}

\subsection{Threshold: Otsu and a fixed floor}
\autofig{diagrams}{F1_threshold_histogram}
Otsu's method picks the level $t$ that splits the 256-bin histogram into a dark and a bright class by maximising the between-class variance
$\sigma_b^2(t)=\omega_0(t)\,\omega_1(t)\,[\mu_0(t)-\mu_1(t)]^2$ ($\omega$ class probabilities, $\mu$ class means). The threshold is $\max(204,\ t_\text{Otsu}+1)$: the floor (heat 0.80) stops a dark frame from producing targets,
Otsu raises the level in warm clutter (T132: bimodal image, $t=40$).
\lstOtsu

\subsection{Despeckle}
A hot pixel with fewer than two hot 8-neighbours is removed: single-pixel noise disappears, real blobs survive (T130: 396 speckle pixels removed, three discs kept).
\lstDespeckle

\subsection{Connected components: two-pass union-find}
\autofig{diagrams}{F3_union_find}
The \emph{first pass} scans row by row; a hot pixel takes the smallest root among its labelled 8-neighbours (W, NW, N, NE) or a new label, and
\emph{records those labels as equivalent} in a union-find forest (path compression; the larger root points to the smaller). The \emph{second pass} replaces each label by its root and renumbers $1..n$. A U-shaped blob whose arms got different labels is merged (T131).
\lstLabel

\subsection{Blob features: moments, orientation, elongation}
\autofig{diagrams}{F4_blob_moments}
One pass accumulates $n,\sum x,\sum y,\sum x^2,\sum y^2,\sum xy$ per component. The centroid is $(\bar x,\bar y)=(\sum x/n,\sum y/n)$ and the second central moments
$\mu_{20}=\frac{\sum x^2}{n}-\bar x^2$, $\mu_{02}=\frac{\sum y^2}{n}-\bar y^2$, $\mu_{11}=\frac{\sum xy}{n}-\bar x\bar y$ form the covariance matrix $C=\begin{bmatrix}\mu_{20}&\mu_{11}\\\mu_{11}&\mu_{02}\end{bmatrix}$ whose eigenvalues are
\begin{equation}\lambda_{1,2}=\frac{\mu_{20}+\mu_{02}}{2}\pm\sqrt{\Bigl(\frac{\mu_{20}-\mu_{02}}{2}\Bigr)^2+\mu_{11}^2}.\end{equation}
The \emph{orientation} of the major axis is $\tfrac12\operatorname{atan2}(2\mu_{11},\mu_{20}-\mu_{02})$ and the \emph{elongation} is $\sqrt{\lambda_1/\lambda_2}$ ($\lambda_2$ is floored at $1/12$, the variance of a single pixel). A blob is a target if its area is 2--2500 pixels (bigger is a fire), it contains a pixel $\ge230$ (an engine) and, from 6 pixels up, its elongation is at most 6 (thin streaks are tracers). T132 checks a bar: orientation $0^\circ$, elongation 10.3. The normalised coordinates
$u=2(\bar x+0.5)/320-1$, $v=1-2(\bar y+0.5)/240$ go to the geometry step; Listing~\ref{lst:Detect} is the whole image stage.
\lstDetect

\subsection{From pixel to ground point}
\sheetfig{s18_params_footprint}
The camera looks along the pitch node's $+Z$. A pixel $(u,v)$ defines the ray direction $\vect d=\mathrm{normalize}\bigl(W_\text{cam}\,(-u\,t_x,\ v\,t_y,\ 1)^\top\bigr)$ with $t_y=\tan(20^\circ)$, $t_x=t_y\cdot320/240$ (image right is local $-X$). The ground point
is where the ray meets the terrain, found by marching in 5\,cm steps until $y\le h(x,z)+h_b$ (Listing~\ref{lst:RayGround}). $h_b=0.2$\,m (\code{kSquadBlobHeight}) is the height of the rover's hot body; stopping at the ground itself would put a single-view fix about 0.27\,m beyond the rover. The standard deviation grows with range, $\sigma=0.12+0.015\,r$\,m.
\lstRayGround

\subsection{Kalman filter: one track per rover}
Each track holds $\vect x=(x,z,v_x,v_z)^\top$ and a $4\times4$ covariance $P$. The prediction over $\Delta t$ uses a constant-velocity model with white-noise acceleration of standard deviation $\sigma_a=1.5$\,m/s$^2$:
\begin{equation}
\vect x^-=F\vect x,\quad P^-=FPF^\top+Q,\quad F=\begin{bmatrix}I_2&\Delta t\,I_2\\0&I_2\end{bmatrix},\quad
Q=\sigma_a^2\begin{bmatrix}\tfrac{\Delta t^4}{4}I_2&\tfrac{\Delta t^3}{2}I_2\\\tfrac{\Delta t^3}{2}I_2&\Delta t^2I_2\end{bmatrix}.
\end{equation}
Only the position is measured: $H=[I_2\ 0]$, $R=\sigma^2I_2$. The update is
\begin{equation}
\vect y=\vect z-H\vect x^-,\quad S=HP^-H^\top+R,\quad K=P^-H^\top S^{-1},\quad \vect x=\vect x^-+K\vect y,\quad P=(I-KH)P^-.
\end{equation}
$S$ is $2\times2$ and inverted in closed form. A new track starts at the measurement with position variance $\sigma^2$ and velocity variance $4\,\text{m}^2/\text{s}^2$. T135: position error below 0.25\,m (late maximum 0.164\,m) and velocity 0.96\,m/s for a
1.2\,m/s target; the estimated velocity lets the squad aim ahead (Chapter~\ref{sec:squad}).
\lstPredict
\lstKalman
\begin{figure}[h]
\centering
\begin{subfigure}{0.49\linewidth}\picdia{kalman_demo}{\linewidth}\caption{noisy measurements, filtered track, velocity}\end{subfigure}
\begin{subfigure}{0.49\linewidth}\picdia{gating_ellipse}{\linewidth}\caption{gating ellipse and measurements}\end{subfigure}
\caption{The Kalman tracker: raw ground points (dots), filtered track (line), $\chi^2$ gate (ellipse).}
\end{figure}

\subsection{Data association: Mahalanobis gating and the Hungarian algorithm}
A measurement $\vect z$ may belong to a track only if the squared Mahalanobis distance of its innovation, $d^2=\vect y^\top S^{-1}\vect y$, is below $\chi^2_{2,0.99}=9.21$ (two degrees of freedom; 99\,\% of true measurements pass). The gate scales with the track's uncertainty (the ellipse above). For each sensor the gated costs $d^2$ form a track-by-measurement matrix and the \textbf{Hungarian algorithm}
(Kuhn--Munkres, $O(n^3)$, padded square) finds the one-to-one assignment of minimum total cost, so two nearby blobs cannot claim one track (T133: optimal in 50 of 50 random matrices against brute force over 120 permutations).
\lstMahal
\lstHungarian

\subsection{Two-view triangulation}
When both UAVs measure the same track in the same frame, the pixel rays $\vect o_1+s\vect d_1$ and $\vect o_2+t\vect d_2$ are combined. Skew lines rarely intersect, so the point is the midpoint of closest approach: minimising $|\vect o_1+s\vect d_1-\vect o_2-t\vect d_2|^2$ with
$\vect w=\vect o_1-\vect o_2$, $a=\vect d_1\!\cdot\!\vect d_1$, $b=\vect d_1\!\cdot\!\vect d_2$, $c=\vect d_2\!\cdot\!\vect d_2$, $d=\vect d_1\!\cdot\!\vect w$, $e=\vect d_2\!\cdot\!\vect w$ gives the $2\times2$ system
\begin{equation}
\begin{bmatrix}a&-b\\b&-c\end{bmatrix}\begin{bmatrix}s\\t\end{bmatrix}=\begin{bmatrix}-d\\-e\end{bmatrix}
\ \Rightarrow\ s=\frac{be-cd}{ac-b^2},\quad t=\frac{ae-bd}{ac-b^2},
\end{equation}
and $\vect p=\tfrac12(\vect o_1+s\vect d_1+\vect o_2+t\vect d_2)$ with the \emph{residual} $|\vect p_1-\vect p_2|$ as a quality measure. Parallel rays ($ac-b^2\approx0$) and points behind a camera ($s$ or $t<0$) are rejected.
The fused measurement has $\sigma=\max(0.08,\,0.5\min(\sigma_1,\sigma_2)+0.25\,\text{residual})$, tighter than a single ray (T134: 0.051\,m against 0.071\,m); orthogonal rays give the best geometry, hence the wing UAV's $90^\circ$ offset.
\lstTriangulate
\begin{figure}[h]
\centering
\picdia{triangulation}{0.75\linewidth}
\caption{Two rays from two UAVs, their closest points, the midpoint and the residual; the ground-ray estimate for comparison.}
\end{figure}

\begin{figure}[h]
\centering
\picsheet{sheet_cv}{\linewidth}
\caption{\sheetcaption{sheet_cv}}
\end{figure}
\subsection{Track life cycle and merging}
\autofig{figures_lab}{cv_strip}
\sheetfig{s18_tests}
An unmatched measurement starts a \emph{tentative} track; after 3 hits it is \emph{confirmed}; with no update for 0.25\,s it \emph{coasts} (it keeps predicting; the HUD shows it amber with a ``?'') and is deleted after 3\,s. A tentative track that misses 2 frames is deleted at once (T136). Two tracks closer than 1\,m with velocities within 1\,m/s are one target: the younger merges into the older, which keeps its id (the squad's lock is on that id). A tentative younger track is fused into the older as one more Kalman measurement, so the converged velocity survives; a confirmed one hands over its whole state. Ids are never reused.
\begin{figure}[h]
\centering
\begin{subfigure}{0.49\linewidth}\piclab{cv_tracks}{\linewidth}\caption{tracks over the thermal image of UAV 0}\end{subfigure}
\begin{subfigure}{0.49\linewidth}\picfive{thermal_fpv_tracks}{\linewidth}\caption{game capture with all track boxes}\end{subfigure}
\caption{Tracks as the player sees them: boxes with ids, range and ``2-VIEW'' when triangulated.}
\end{figure}
